\newif\ifuseseminar
\useseminartrue
\input{../../latex_common/header.tex.frag}

\title{Relatividade Restrita -- Lista de Exercícios: Dilatação do Tempo, Contração do Espaço, Paradoxo dos Gêmeos e Efeito Doppler}

\begin{document}

Em toda a lista abaixo usaremos assinaturas $(-,+,+,+)$ e as coordenadas em anos-luz.

\begin{enumerate}

      \item Quadrivelocidade e tempo próprio: Seja $x^\mu(\tau)$ a trajetória de uma
            partícula parametrizada pelo seu tempo próprio $\tau$ e defina a
            quadrivelocidade por
            $$
                  u^\mu \equiv \frac{1}{c}\frac{\mathrm{d}x^\mu}{\mathrm{d}\tau}.
            $$

            \begin{enumerate}
                  \item Mostre que $u^\mu$ transforma-se como um quadrivetor sob
                        transformações de Lorentz.

                  \item Mostre que
                        $$
                              u^\mu=(\gamma,\gamma\beta^1,\gamma\beta^2,\gamma\beta^3),
                        $$
                        onde $\beta^i=v^i/c$.

                  \item Utilizando a métrica de Minkowski, demonstre que
                        $$
                              u^\mu u_\mu=-1.
                        $$

                  \item Mostre que, para uma partícula em repouso,
                        $$
                              u^\mu=(1,0,0,0).
                        $$
            \end{enumerate}

      \item Quadriaceleração: Defina a quadriaceleração por
            $$
                  a^\mu\equiv\frac{\mathrm d u^\mu}{\mathrm d\tau}.
            $$

            \begin{enumerate}
                  \item Mostre que $a^\mu$ transforma-se como um quadrivetor.

                  \item Derive a relação
                        $$
                              u^\mu a_\mu=0.
                        $$

                  \item Interprete geometricamente o resultado obtido.
            \end{enumerate}

      \item Quadrimomento: Defina o quadrimomento por
            $$
                  p^\mu\equiv mc\,u^\mu.
            $$

            \begin{enumerate}
                  \item Justifique fisicamente essa definição.

                  \item Mostre que
                        $$
                              p^\mu=
                              \left(
                              \frac{E}{c},
                              \mathbf p
                              \right).
                        $$

                  \item Demonstre que
                        $$
                              p^\mu p_\mu=-m^2c^2.
                        $$

                  \item Mostre que essa quantidade é invariante sob transformações de
                        Lorentz.
            \end{enumerate}

      \item Relação energia-momento:

            \begin{enumerate}
                  \item Utilizando o resultado do exercício anterior, demonstre que
                        $$
                              E^2=p^2c^2+m^2c^4.
                        $$

                  \item Obtenha o limite não relativístico dessa expressão.

                  \item Considere o caso $m=0$.

                  \item Mostre que, nesse caso,
                        $$
                              E=pc.
                        $$
            \end{enumerate}

      \item Energia relativística:

            \begin{enumerate}
                  \item Considere uma partícula de massa $m$ e velocidade $\mathbf{v}$.
                        Utilizando a definição do quadrimomento
                        $$
                              p^\mu = mc\,u^\mu,
                        $$
                        mostre que sua componente temporal pode ser escrita como
                        $$
                              p^0 = \gamma mc.
                        $$

                  \item Definindo a energia relativística por
                        $$
                              E \equiv cp^0,
                        $$
                        demonstre que
                        $$
                              E = \gamma mc^2.
                        $$

                  \item Utilizando a relação
                        $$
                              p = \gamma mv,
                        $$
                        elimine $\gamma$ e obtenha a expressão
                        $$
                              E^2 = p^2c^2 + m^2c^4.
                        $$

                  \item Mostre que, no limite $v \ll c$,
                        $$
                              E \simeq mc^2 + \frac12 mv^2.
                        $$
                        Interprete fisicamente os dois termos desta expressão.
            \end{enumerate}

      \item Quadriforça: Defina a quadriforça por
            $$
                  F^\mu\equiv\frac{\mathrm d p^\mu}{\mathrm d\tau}.
            $$

            \begin{enumerate}
                  \item Mostre que $F^\mu$ transforma-se como um quadrivetor.

                  \item Demonstre que
                        $$
                              F^\mu u_\mu=0.
                        $$

                  \item Relacione as componentes de $F^\mu$ com a força
                        tridimensional.
            \end{enumerate}

      \item Quadrivetor de onda e efeito Doppler: Considere uma onda plana cuja fase é
            $$
                  \phi=k_\mu x^\mu.
            $$

            \begin{enumerate}
                  \item Mostre que a invariância da fase implica que $k^\mu$ é um
                        quadrivetor.

                  \item Escreva
                        $$
                              k^\mu=
                              \left(
                              \frac{\omega}{c},
                              \mathbf k
                              \right).
                        $$

                  \item Obtenha a fórmula relativística para o efeito Doppler a partir
                        de uma transformação de Lorentz.
            \end{enumerate}

      \item Fótons:

            \begin{enumerate}
                  \item Mostre que, para uma onda eletromagnética no vácuo,
                        $$
                              k^\mu k_\mu=0.
                        $$

                  \item Utilizando
                        $$
                              p^\mu=\hbar k^\mu,
                        $$
                        mostre que
                        $$
                              p^\mu p_\mu=0.
                        $$

                  \item Demonstre que
                        $$
                              E=pc.
                        $$

                  \item Discuta a interpretação geométrica desse resultado.
            \end{enumerate}

      \item Quadripotencial eletromagnético: Defina o quadripotencial
            $$
                  A^\mu=
                  \left(
                  \frac{\phi}{c},
                  \mathbf A
                  \right).
            $$

            \begin{enumerate}
                  \item Mostre que $A^\mu$ transforma-se como um quadrivetor.

                  \item Identifique suas componentes temporal e espaciais.

                  \item Discuta a liberdade de calibre.
            \end{enumerate}

      \item Tensor eletromagnético: Defina
            $$
                  F_{\mu\nu}
                  =
                  \partial_\mu A_\nu
                  -
                  \partial_\nu A_\mu.
            $$

            \begin{enumerate}
                  \item Mostre que $F_{\mu\nu}$ é antissimétrico.

                  \item Escreva explicitamente sua matriz de componentes em termos dos
                        campos elétrico e magnético.

                  \item Mostre que $F_{\mu\nu}$ transforma-se tensorialmente sob
                        transformações de Lorentz.
            \end{enumerate}

      \item Transformações dos campos eletromagnéticos:

            \begin{enumerate}
                  \item Considere um boost ao longo do eixo $x$.

                  \item Utilize a transformação tensorial de $F_{\mu\nu}$ para obter
                        as transformações dos campos elétrico e magnético.

                  \item Mostre que componentes de $\mathbf E$ e $\mathbf B$
                        misturam-se sob transformações de Lorentz.

                  \item Interprete esse resultado fisicamente.
            \end{enumerate}

      \item Invariantes do campo eletromagnético:

            \begin{enumerate}
                  \item Mostre que
                        $$
                              F_{\mu\nu}F^{\mu\nu}
                        $$
                        é um invariante de Lorentz.

                  \item Escreva esse invariante em termos de $\mathbf E$ e
                        $\mathbf B$.

                  \item Mostre que
                        $$
                              \epsilon^{\mu\nu\rho\sigma}
                              F_{\mu\nu}
                              F_{\rho\sigma}
                        $$
                        também é um invariante de Lorentz.

                  \item Escreva esse segundo invariante em termos dos campos
                        elétrico e magnético.
            \end{enumerate}

      \item Equações de Maxwell com fontes: Defina a quadricorrente
            $$
                  J^\mu=(c\rho,\mathbf J).
            $$

            \begin{enumerate}
                  \item Mostre que as equações de Maxwell com fontes podem ser
                        escritas como
                        $$
                              \partial_\mu F^{\mu\nu}
                              =
                              \mu_0J^\nu.
                        $$

                  \item Recupere a lei de Gauss.

                  \item Recupere a lei de Ampère-Maxwell.
            \end{enumerate}

      \item Equações homogêneas de Maxwell: Defina o tensor dual
            $$
                  \widetilde F^{\mu\nu}
                  =
                  \frac12
                  \epsilon^{\mu\nu\rho\sigma}
                  F_{\rho\sigma}.
            $$

            \begin{enumerate}
                  \item Mostre que
                        $$
                              \partial_\mu\widetilde F^{\mu\nu}=0.
                        $$

                  \item Recupere a lei de Gauss para o magnetismo.

                  \item Recupere a lei de Faraday.
            \end{enumerate}

\end{enumerate}

\end{document}
